\section{Localized alias inverse theorems}
\label{sec:tpc62-localized-alias}

TPC--61 identifies the exact normalized missing Gram, but deliberately does
not estimate its off-diagonal arithmetic kernel \citep{WangTPC61}.  This
section records the converse information that is available without any new
number theory: if the restricted native norm is large, then a large portion
of that norm must already be visible on one row, on one literal collision
cell, and finally on one row pair.  Thus a failure of a proposed norm bound
cannot remain completely delocalized.

Let \(F_M\) be the normalized physical missing incidence matrix on the
active row space, and put
\begin{equation}
 \mathbf K_M:=F_M^*F_M,
 \qquad
 \mathbf H_M:=\mathbf K_M-I,
 \qquad
 \kappa_M:=\lambda_{\max}(\mathbf K_M).
 \label{eq:tpc62-normalized-alias-gram}
\end{equation}
Every column of \(F_M\) has norm one.  Hence
\begin{equation}
 (\mathbf K_M)_{m,m}=1,
 \qquad
 (\mathbf H_M)_{m,m}=0,
 \qquad
 \kappa_M=1+\lambda_{\max}(\mathbf H_M).
 \label{eq:tpc62-zero-diagonal-alias}
\end{equation}
The rows in this section are only the actually active rows.  No zero row is
adjoined to inflate a degree or a dimension.

\subsection{Literal collision-cell decompositions}

We use a decomposition that is exact at the level of Gram summands.  Let
\(\mathscr B\) be a nonempty set of \(B\) literal collision cells and write
\begin{equation}
 \mathbf H_M=\sum_{b\in\mathscr B}\mathbf H_M^{(b)}.
 \label{eq:tpc62-cell-decomposition}
\end{equation}
Here every common-output product of two distinct source rows is assigned to
exactly one cell, together with its conjugate product, and
\((\mathbf H_M^{(b)})_{m,m}=0\).  Thus
\eqref{eq:tpc62-cell-decomposition} is an identity with the literal phases;
it is not a support majorant.  If the original packet is first partitioned
into source cells, then \(\mathscr B\) must include all nonempty cell pairs
that can contribute a Gram product.  In particular, \(B\) below counts
correlation pieces, not merely the number of source-cell names.  The
decomposition is fixed by the literal carrier before any coefficient vector
or extremizing eigenvector is selected.

For a row \(m\) and a cell \(b\), set
\begin{equation}
 \mathcal N_b(m)
 :=\{n\ne m:(\mathbf H_M^{(b)})_{m,n}\ne0\},
 \qquad
 \Delta:=\max_{m,b}|\mathcal N_b(m)|.
 \label{eq:tpc62-cell-neighbour-degree}
\end{equation}
This is the degree of the actual cellwise collision graph.  It is not an
ambient number of labels that never occur.

\begin{theorem}[Localized alias inverse theorem]
\label{thm:tpc62-localized-alias-inverse}
Assume the active row space is nonzero and
\begin{equation}
 \kappa_M\ge 1+L
 \label{eq:tpc62-large-kappa-hypothesis}
\end{equation}
for some \(L>0\).  Then all of the following hold.
\begin{enumerate}[label=\textup{(\roman*)},leftmargin=3em]
 \item There is an active row \(m\) such that
 \begin{equation}
  \sum_{n\ne m}|(\mathbf K_M)_{m,n}|\ge L.
  \label{eq:tpc62-large-row-l1}
 \end{equation}
 \item There are an active row \(m\) and a literal cell
       \(b\in\mathscr B\) such that
 \begin{equation}
  \sum_{n\ne m}|(\mathbf H_M^{(b)})_{m,n}|
  \ge \frac{L}{B}.
  \label{eq:tpc62-large-cell-row-l1}
 \end{equation}
 \item If \(\Delta<\infty\), then for some \(m\ne n\) and some
       \(b\in\mathscr B\),
 \begin{align}
  |(\mathbf H_M^{(b)})_{m,n}|
  &\ge \frac{L}{B\Delta},
  \label{eq:tpc62-large-pair-witness}\\
  \left(\sum_{n\in\mathcal N_b(m)}
  |(\mathbf H_M^{(b)})_{m,n}|^2\right)^{1/2}
  &\ge \frac{L}{B\sqrt{\Delta}}.
  \label{eq:tpc62-large-l2-star}
 \end{align}
\end{enumerate}
When \(L>0\), the conclusion itself forces \(\Delta\ge1\).
\end{theorem}

\begin{proof}
The Hermitian matrix \(\mathbf H_M\) has zero diagonal.  Its largest
eigenvalue is at most its spectral radius, and the spectral radius is at
most the induced \(\ell^\infty\)-norm.  Therefore
\[
 L
 \le\lambda_{\max}(\mathbf H_M)
 \le\max_m\sum_{n\ne m}|(\mathbf H_M)_{m,n}|.
\]
This proves \eqref{eq:tpc62-large-row-l1}, because the off-diagonal entries
of \(\mathbf H_M\) and \(\mathbf K_M\) agree.

Choose a row \(m\) satisfying \eqref{eq:tpc62-large-row-l1}.  The exact
cell decomposition and the triangle inequality give
\[
 L
 \le\sum_{n\ne m}\left|
       \sum_{b\in\mathscr B}(\mathbf H_M^{(b)})_{m,n}\right|
 \le\sum_{b\in\mathscr B}\sum_{n\ne m}
       |(\mathbf H_M^{(b)})_{m,n}|.
\]
At least one cell has the mass in
\eqref{eq:tpc62-large-cell-row-l1}.  That cell has at most \(\Delta\)
nonzero neighbours, so its largest entry is at least its \(\ell^1\)-mass
divided by \(\Delta\), proving
\eqref{eq:tpc62-large-pair-witness}.  Finally Cauchy--Schwarz on the same
star gives
\[
 \frac{L}{B}
 \le\sum_{n\in\mathcal N_b(m)}
       |(\mathbf H_M^{(b)})_{m,n}|
 \le\sqrt{\Delta}
 \left(\sum_{n\in\mathcal N_b(m)}
       |(\mathbf H_M^{(b)})_{m,n}|^2\right)^{1/2},
\]
which is \eqref{eq:tpc62-large-l2-star}.
\end{proof}

The theorem is one-sided.  A large cell entry is compatible with
cancellation between cells and therefore need not force a comparably large
entry of the completed Gram.  What the theorem says is that a large
\emph{completed} norm forces at least one large literal cell sum.

\subsection{Degree times pair correlation}

The forward Schur estimate is the exact counterpart of the inverse theorem.

\begin{corollary}[Pair correlation plus active degree]
\label{cor:tpc62-degree-pair-exponent}
Suppose, uniformly over all colliding distinct active rows, that
\begin{align}
 |\{n\ne m:(\mathbf K_M)_{m,n}\ne0\}|
 &\le X^{\delta+o(1)},
 \label{eq:tpc62-global-degree-exponent}\\
 |(\mathbf K_M)_{m,n}|
 &\le X^{-\tau+o(1)},
 \label{eq:tpc62-pair-correlation-exponent}
\end{align}
where the two \(o(1)\) terms are uniform in the row indices.  Then
\begin{equation}
 \boxed{
 \kappa_M
 \le 1+X^{\delta-\tau+o(1)},
 \qquad
 \nu_M\le(\delta-\tau)_+,}
 \label{eq:tpc62-degree-pair-kappa-bound}
\end{equation}
where
\(
 \nu_M:=\limsup_{X\to\infty}
 \log\max\{1,\kappa_M\}/\log X
\).
In particular, a pair saving with \(\tau<\delta\) does not by itself give
a subpower operator bound: this Schur certificate still permits the
positive exponent \(\delta-\tau\).  It does not assert that the actual
\(\nu_M\) is positive.
\end{corollary}

\begin{proof}
Every off-diagonal absolute row sum is at most
\(X^{\delta-\tau+o(1)}\).  The same induced-norm estimate used in
\cref{thm:tpc62-localized-alias-inverse} yields the first assertion.  If
\(\delta\le\tau\), the right side is \(X^{o(1)}\); otherwise its exponent
is \(\delta-\tau\).  This proves the second assertion.
\end{proof}

If instead the available estimate is cellwise, with
\(B=X^{b+o(1)}\), cell degree \(X^{\delta+o(1)}\), and
\(|(\mathbf H_M^{(b)})_{m,n}|\le X^{-\tau+o(1)}\), then the same proof
gives
\begin{equation}
 \nu_M\le(b+\delta-\tau)_+.
 \label{eq:tpc62-cellwise-exponent-bound}
\end{equation}
The factor \(B\) may be omitted only when orthogonality or another exact
summation mechanism removes it.

Conversely, suppose
\begin{equation}
 \kappa_M\ge1+X^{\nu-o(1)},
 \qquad
 B\le X^{b+o(1)},
 \qquad
 \Delta\le X^{\delta+o(1)}.
 \label{eq:tpc62-exponent-inverse-hypotheses}
\end{equation}
Then \cref{thm:tpc62-localized-alias-inverse} produces one cell and one
row pair for which
\begin{align}
 |(\mathbf H_M^{(b)})_{m,n}|
 &\ge X^{\nu-b-\delta-o(1)},
 \label{eq:tpc62-exponent-pair-witness}\\
 \left(\sum_{n\in\mathcal N_b(m)}
 |(\mathbf H_M^{(b)})_{m,n}|^2\right)^{1/2}
 &\ge X^{\nu-b-\delta/2-o(1)}.
 \label{eq:tpc62-exponent-star-witness}
\end{align}
These are deterministic localization statements; no probability model is
being invoked.

\subsection{The arithmetic content of the localized witness}

Conjugating \(\mathbf H_M\) and every cell piece by the same fixed diagonal
source-divisor gauge preserves the decomposition, every entry magnitude,
and all row/star norms above.  After this TPC--61 gauge, each direct-slice
contribution to an off-diagonal entry is exactly the normalized,
fixed-\(h_0\) two-affine M\"obius sum
\begin{equation}
 \frac{1}{\sqrt{W_M(m)W_M(m')}}
 \sum_j
 \mathfrak m_m(j)\mathfrak m_{m'}(j)
 \mu(mj+h_0)\mu(m'j+h_0)
 \overline{a_m(j)}a_{m'}(j)\Gamma_{m,m'}(j).
 \label{eq:tpc62-fixed-shift-pair-kernel}
\end{equation}
The full off-diagonal entry is the sum of the left and right direct-slice
expressions on their declared orthogonal output components; no cross-slice
Gram term is inserted.
Equivalently, the same sum is indexed by the literal solutions of
\begin{equation}
 m'pr-mp'r'=h_0(m'-m),
 \label{eq:tpc62-determinant-surface}
\end{equation}
with all inherited largest-prime, cofactor, missing-carrier, orbit, and
endpoint restrictions retained.  Restricting the Gram summands to a
literal collision cell gives the entries
\((\mathbf H_M^{(b)})_{m,m'}\) above.  Consequently
\eqref{eq:tpc62-exponent-pair-witness} localizes a large restricted native
norm to one normalized signed determinant sum.  It does not localize to one
determinant quadruple, because many quadruples may still occur in that sum.

This is an exact reformulation, not a cancellation theorem.  TPC--61 does
not prove \eqref{eq:tpc62-pair-correlation-exponent} for any
\(\tau>0\), either for one pair uniformly or for the growing matrix of row
pairs.  The present inverse theorem likewise supplies no sign cancellation:
it only says where a large norm must leave a witness.  The matrix statements
are L0, their literal fixed-shift identification is L1, and no L2 parity or
prime-pair improvement is claimed.
